\section{Three connected research questions and verified gap}
Commercial drone-delivery operators sharing capacity-constrained urban low-altitude corridors can avoid delay by sharing flight intent, yet route and timing detail may reveal demand to competitors. Operators, the traffic-management authority, customers, nearby/public stakeholders, and emergency/public-safety users are affected. Closest verified work designs efficient/fair traffic-flow protocols, limited-information traffic management, privacy-preserving coordination, and vertiport auctions \cite{chin2022efficient,chin2023traffic,maheshwari2025privacy,su2024vertiport}; oligopoly research explains incentives to share private information \cite{raith1996general}. The narrow surviving gap is endogenous \emph{granularity} of voluntary, non-safety disclosure when commercial-confidentiality cost is private and coordination value changes with congestion---not a claim that privacy or airspace auctions are new.

\noindent\textbf{Q1---Economics:} Under incomplete information about commercial-confidentiality costs, how does congestion change equilibrium disclosure, and when does equilibrium disclosure differ from the full-information first best?\par
\noindent\textbf{Q2---Computation:} How do equilibrium disclosure, welfare loss, and candidate-mechanism performance vary across congestion, confidentiality cost, and access-rule parameters?\par
\noindent\textbf{Q3---Behavior:} Do participants depart from the Bayesian benchmark because of privacy salience, reciprocity, fairness, competitive concerns, or misunderstanding of strategic uncertainty?\par
\noindent\textbf{Integrated question:} Which disclosure-and-access design remains welfare-improving, behaviorally intelligible, and institutionally accountable across private costs and congestion? Figure~\ref{fig:teaser} connects the questions: incentives define the model, computation tests it, behavior probes its assumptions, and allocation plus a revision gate bound the practical claim.

\section{Answer to Q1: incentives and strategic model}
Players $i\in\{A,B\}$ simultaneously choose $a_i\in\{M,P,F\}$ (Minimum, Partial, Full voluntary commercial disclosure); mandatory safety information is outside the action set. Nature independently draws private $c_i\in\{1,1.6\}$ with prior $(0.5,0.5)$; public congestion $d\in\{0.3,0.6,1\}$ is known before choice. A complete pure strategy maps each own type to an action. With $x=(0,1,2)$, $k=(0,1,2.5)$, $B=7.5$, and $\lambda=0.5$,
\[
u_i(a_i,a_j;c_i,d)=dB[1-e^{-\lambda(x_i+x_j)}]-c_i k(a_i).
\]
Simultaneous choice plus an unobserved rival confidentiality type makes pure-strategy Bayesian Nash equilibrium (BNE) the appropriate benchmark \cite{harsanyi1967}. Exhaustive interim best-response enumeration gives the unique type-contingent BNE $(L,H)$: Low $(M,M)$, Medium $(P,M)$, High $(P,P)$.

\textbf{Game theory} predicts $M/P/F$ disclosure by type. \textbf{Social choice} compares $W=u_A+u_B$ with the realized-type full-information first best, exposing coordination/privacy and efficiency/fairness trade-offs: BNE/first-best/gap is $0/0.639/0.639$, $2.193/3.139/0.946$, and $6.882/7.385/0.503$. \textbf{Mechanism design} changes enhanced non-safety information access through CCAR and tests incentive/welfare robustness: at Medium, $\alpha_M=0.7$ changes BNE to $(P,P)$ and underlying welfare to $3.089$ (gap $0.050$), without restricting safety information.

\section{Answer to Q2: computation, auction comparison, and artifacts}
\textbf{Language-independent pseudocode:} for each congestion state, enumerate type-contingent strategies; compute interim utility over rival types; retain mutual best responses as pure BNE; for each realized type pair enumerate action profiles and maximize full-information welfare; compare expected BNE and first-best welfare; repeat under CCAR and parameter sweeps. Inputs are the stated primitives; baseline is unrestricted enhanced access, comparison is CCAR, and metrics are equilibrium actions, welfare/gap, multiplicity, and failure diagnostics. Fresh runs yield the results above; a 760-cell sweep finds 401 no-effect, 334 robust-improvement, 18 excessive-disclosure, 3 increased-gap, 66 multiple-BNE, and 0 no-pure-BNE cells (overlapping diagnostics, not probabilities).

The bounded application allocates one residual commercial priority slot (emergency/public-safety users excluded) among operators with private values and sealed reports. FCFS allocates to the first eligible request and stops; a single-unit second-price auction allocates to the highest bid, charges the second bid, breaks ties by ID, and stops. Under the independent-private-values rational benchmark, truthful bidding is weakly dominant \cite{vickrey1961counterspeculation}. For values $(8,5,3,2)$ and arrival $B,C,A,D$, FCFS gives $B$ efficiency $0.625$; truthful second price gives $A$, payment $5$, utility $3$, efficiency $1$. In 10,000 seed-206 rounds, mean efficiency is $0.622308$ versus $1.000000$ and auction revenue is $5.99207$. Ability to pay, urgency, interdependent values, and legitimacy bound the analogy.

\textbf{Artifacts:} \PublicRepo{} at current reproducibility checkpoint \CurrentCommit; \ColabOne; \SpaceRepo; validated editable poster \texttt{\seqsplit{PS2-FP10-Yiqiao-A0-Poster.pptx}}. The two notebook paths, both hosted Run-all confirmations, and all versions appear in Appendix~\ref{app:technical}.

\balance\small
\section{Answer to Q3: behavior and interdisciplinary integration}
The rational benchmark is type-contingent BNE disclosure. Candidate mechanisms are \emph{privacy salience}, \emph{reciprocity/fairness}, and \emph{competitive concern}; the competing explanation is misunderstanding incomplete-information incentives or the payoff structure. The Hugging Face Solo and same-device Peer Play environment is exploratory evidence infrastructure: choice precedes benchmark reveal, then CCAR and reflection probe responses. No representative behavioral sample exists. A planned discriminating test compares initial choice, post-benchmark reflection, CCAR second choice, and a trusted-partner counterfactual: stable deviations after comprehension checks that move with framing/partner trust favor preference or framing accounts; convergence after explanation favors misunderstanding. Evidence would qualify the benchmark or revise the interface, not establish population effects.

\section{Advanced development, real-world impact, and roadmap}
For commercial drone-delivery operators sharing capacity-constrained urban low-altitude corridors, the decision is how much non-safety flight-intent information to share and how to allocate residual commercial priority scarcity. Potential benefit is coordination and avoided delay; risks are commercial exposure, unequal access, over-disclosure, and misspecified incentives. Before adoption, the authority, operators, customers, nearby/public stakeholders, and emergency/public-safety users need pilot calibration of latency and confidentiality value, safety review, and distributional analysis. Correction requires auditable rules, a safety-information floor, appeal/review, and parameter revision or suspension when welfare/access diagnostics fail. Academia requires equilibrium, identification, robustness, and reproducibility; industry requires feasibility, reliability, confidentiality, and latency; public institutions require safety, access, legitimacy, auditability, and correction. Next: preregister the behavioral discrimination above and pilot the mechanism against baseline access.

\noindent\textbf{Expected 2056 Nobel Prize and Turing Award contribution statement (aspirational).} Building on Harsanyi's incomplete-information foundation \cite{harsanyi1967} and Vickrey's incentive-compatible allocation \cite{vickrey1961counterspeculation}, we aspire by 2056 to develop auditable adaptive institutions that join private-information economics, reproducible computation, and behavioral diagnosis. Recognition would require general theory, independent replication, field evidence, equitable safety outcomes, and durable public correction mechanisms; none is claimed achieved here.

\subsection*{Open Science Statement}
\PublicRepo{}, \ColabOne, \SpaceRepo{}, and \texttt{\seqsplit{PS2-FP10-Yiqiao-A0-Poster.pptx}} point to the same stylized/synthetic model. Checkpoint \CurrentCommit{} uses Python 3.13.7, NumPy 2.3.5, Matplotlib 3.10.8, nbclient 0.10.2, and Gradio 6.29.0. Shortest reproduction: open \ColabOne{} and Run all; outputs include BNE, welfare, CCAR, sensitivity, and auction tables/figures. Limits include uncalibrated parameters, two players/types, pure strategies, and no representative behavioral sample.

\subsection*{Statement of Contribution to the UN's SDGs}
No demonstrated SDG contribution is claimed at proposal stage. Prospective relevance to safe, inclusive urban infrastructure would require a specified target, affected audience, access indicator, safety and confidentiality audit, distributional baseline, field comparison, and a public plan for unintended effects and suspension.
